\section{An actual SNC block produces its intrinsic old carrier and count}
\label{sec:nc-snc-block-carrier}

\begin{theorem}[The actual union of one SNC block is adapted to the entire other block]
\label{hir:nc-snc-old-block-boundary-adapted-union}
\lean{TranscendentalNakamaye.Pending.HironakaNormalCrossingsSNCBlockCarrier.boundaryAdapted_union_of_snc_append}
\leanok
Let the actual finite concatenation $B+\!+E$ be SNC. Then the literal union
$H=\bigcup_{A\in B}A$ is an NC carrier adapted to the entire retained list
$E$. Empty members and all original index positions are retained. This is
proved from the genuine SNC coordinates, without a joint-union witness.
\uses{hir:nc-original-highest-joint-branch-chart,hir:general-nc-simultaneous-original-boundary}
\end{theorem}
\begin{proof}
The finite union of the closed old members is closed and every retained
member is closed. At an arbitrary point $x$, take the actual SNC flattening
of the whole concatenation. Shrink its actual source so the corresponding
manifold domain avoids every old member not through $x$; the finite
intersection of their open complements is a genuine neighbourhood of $x$.
Restriction of the original local holomorphic equivalence supplies a true
restricted flattening with the same coordinate functions and the same zero
centre.

Let $S$ be the union of the singleton coordinate sets of precisely those old
indices through $x$, and retain the original singleton coordinate set at
every born index. On the restricted domain, the literal union of old members
is exactly the union of the coordinate hyperplanes indexed by $S$: absent
old members have been excluded and every through-old member has its true
coordinate equation. Injectivity on all through-indices makes the singleton
coordinate sets of distinct old and born indices disjoint. In particular
$S$ is disjoint from every through-born singleton, and the born coordinate
sets remain injective. These are all fields of the actual boundary-adapted
carrier, with no reordering or deletion of retained indices.
\end{proof}

\begin{theorem}[Its genuine branch multiplicity is the number of through-old indices]
\label{hir:nc-snc-old-block-intrinsic-count}
\lean{TranscendentalNakamaye.Pending.HironakaNormalCrossingsSNCBlockCarrier.branchMultiplicity_union_eq_count}
\leanok
For the actual carrier just constructed and every actual point $x$,
its intrinsic branch multiplicity is the number of indices $i$ of $B$ with
$x\in B_i$, and consequently is at most the length of $B$.
\uses{hir:nc-snc-old-block-boundary-adapted-union,hir:general-nc-intrinsic-branch-multiplicity}
\end{theorem}
\begin{proof}
Use the same genuine restricted coordinates at $x$. Distinct through-old
indices have distinct singleton coordinate sets, by actual SNC injectivity.
Their finite union $S$ therefore has cardinality equal to the number of
through-old indices; empty and absent members contribute nothing. The true
intrinsic branch-count theorem applied to these centred coordinates gives
multiplicity $|S|$, independent of any chosen flattening. The through-index
set is contained in the whole finite index type of $B$, giving the stated
bound. This counts actual old branches without identifying them with born
members or inferring any later-stage old multiplicity decrease.
\end{proof}

\begin{theorem}[The actual original coefficient phase has genuine current local old carriers]
\label{hir:nc-original-boundary-phase-current-local-adapted-carriers}
\lean{TranscendentalNakamaye.Pending.HironakaNormalCrossingsSNCBlockCarrier.exists_original_phase_current_local_carriers}
\leanok
From only the original natural compact-cosupport positive-mark joint NC data,
one actual finite stratum phase is produced as before. At every original
highest point it internally chooses $U,B$, and at EVERY actual restricted
prefix the strict transform of $H|_U$ is a genuine NC carrier adapted to the
SAME entire current retained/born boundary. Its intrinsic count equals the
number of original old branch indices whose strict transforms pass through
the point and is bounded by the original finite list length.
\uses{hir:nc-original-boundary-stratum-joint-histories,hir:nc-snc-old-block-boundary-adapted-union,hir:nc-snc-old-block-intrinsic-count}
\end{theorem}
\begin{proof}
Internally produce the actual whole-boundary phase and its genuine joint
branch histories. At a chosen original highest point, take their actual
$U,B$. Every restricted prefix is truly well formed. Its literal entire
boundary identity is the strict old list followed by exactly the current
original-and-born list, so its proved SNC property gives SNC of precisely
these two actual blocks. Apply the preceding actual carrier and count
construction to them. The strict transform of a finite union is the union
of the strict transforms: prove this by finite-list induction using the
true binary-union strict-transform identity and the empty case. Together
with $H|_U=\bigcup B$, this identifies the constructed carrier with the
literal current old carrier, without a geometric presentation certificate.
Mapping a finite list preserves its length and indices, giving the count
and bound for exactly the same original old branch slots. This produces
current local geometry but does not assert a global smaller multiplicity
outside the permitted cosupport or complete the open-extension iteration.
\end{proof}
